\resetproblem \begingroup % \jeolmheader
    \def\jeolmdate{сентябрь 2022 г.}%
    \def\jeolmgroupname{9}%
    \def\jeolmauthors{Бибиков~П.\,В.}%
\jeolmheader \endgroup

\worksheet{Комбинаторика}

\begin{problems}

\item Пете и Васе на лбу написали по одной несократимой дроби. Они видят дроби друг друга, но не видят свои дроби. Они должны написать на листе бумаги (каждый --- на своем листе) конечный набор дробей. Могут ли они действовать так, чтобы хотя бы один из них написал дробь, написанную у него на лбу?

\claim{Решение} Пусть у Пети написана несократимая дробь $p_1/q_1$, а у Васи --- дробь $p_2/q_2$. Пусть Петя выпишет на листке бумаги все несократимые дроби $a/b$, где $|a|+|b|\leq|p_2|+|q_2|$, а Вася --- все несократимые дроби $c/d$, где $|c|+|d|\leq|p_1|+|q_1|$. В случае, если $|p_1|+|q_1|\leq |p_2|+|q_2|$, среди выписанных Петей дробей окажется дробь $p_1/q_1$, а в случае, если $|p_1|+|q_1|\geq |p_2|+|q_2|$, дробь $p_2/q_2$ будет выписана у Васи, что и требовалось.

\item В классе 33 ученика, а сумма их возрастов составляет 430 лет. Верно ли, что найдутся 20 учащихся, сумма возрастов которых больше 260 лет?

\claim{Решение} Предположим противное: пусть у любых 20 учеников сумма возрастов меньше 260 лет. Тогда просуммируем возраста всевозможных двадцаток учеников: он будет меньше $C_{33}^{20}\cdot 260$. С дрогой стороны, каждый ученик входит в $C_{32}^{19}$ двадцаток, поэтому сумма возрастов всех двадцаток равна $C_{32}^{19}\cdot 430$. Таким образом, $$C_{33}^{20}\cdot 260>C_{32}^{19}\cdot 430\quad\Longrightarrow\quad 429> 430 $$ --- противоречие.

\item Несколько прямых, никакие две из которых не параллельны, разрезают плоскость на части. Внутри одной из этих частей отметили точку $A$. Докажите, что точка, лежащая с $A$ по разные стороны от всех данных прямых, существует тогда и только тогда, когда часть, содержащая $A$, неограничена.

\claim{Решение} Сначала докажем, что если точка $A$ лежит в ограниченной области, искомой точки для нее не существует. В самом деле, пусть точка $A$ лежит в области, ограниченной прямыми $\ell_1$, \ldots, $\ell_n$. Тогда рассмотрим сначала пару прямых $\ell_1$ и $\ell_2$. Искомая точка должна лежать в угле, образованном этими прямыми. Однако этот угло лежит в той же полуплоскости, что и точка $A$, относительно прямой $\ell_3$. Значит, в этом случае искомой точки не существует.

Теперь докажем, что если точка $A$ лежит в неограниченной области, то искомая точка для нее существует. Для этого посмотрим на конфигурацию прямых издалека. Мы увидим, что все прямые проходят вблизи одной точки, а точка $A$ лежит угле, ограниченном двумя прямыми. Возьмем точку, лежащую в угле, образованном теми же прямыми, но с другой стороны то точки $O$. Легко видеть, что полученная точка будет искомой. 

\item Натуральное число кратно 4. Все его делители выписали в строку в порядке возрастания. Докажите, что в
этой строке найдется пара соседей с разностью 2.

\claim{Решение} Рассмотрим делители 2 и 4. Если наше число не делится на 3, то искомые делители найдены. В противном случае число делится на 3, и мы получаем три подряд идущих делителя: 2, 3 и 4. 

Запустим следующий процесс. Рассмотрим три подряд идущих делителя, средний из которых нечетен. Выберем к нему в пару четный делитель, не кратный 4. Умножим эти делители на 2. Ясно, что новые числа также будут делителями, причем между ними расстояние 2. Если число между ними также является делителем, то мы получаем новую тройку подряд идущих делителей, центральный из которых нечетен. Таким образом, мы увеличиваем тройки делителей, и рано или поздно этот процесс остановится. В этот момент мы и получим пару делителей, отличающихся на 2, что и требовалось.

\item На шахматной доске $n\times n$ расставлены ладьи так, что они бьют все черные клетки. Какое наибольшее возможное количество непобитых белых клеток может быть?

\claim{Решение} Рассмотрим белую непобитую клетку. Поскольку в одной строке с ней есть не молее $\lfloor n/2\rfloor$ черных клеток, каждую из которых должна бить какая-то ладья, то в соответствующих столбцах стоит хотя бы одна ладья. Аналогично, в одном столбце с непобитой белой клеткой есть не более $\lfloor n/2\rfloor$ черных клеток, каждую из которых должна бить какая-то ладья, то в соответствующих строках стоит хотя бы одна ладья. Значит, есть не менее $\lfloor n/2\rfloor$ целиком побитых столбцов и не менее $\lfloor n/2\rfloor$ целиком побитых строк. Значит, останется не более $\lceil n/2\rceil$ целиком непобитых столбцов и не более $\lceil n/2\rceil$ целиком непобитых строк. На их пересечении будет не более $\lceil n/2\rceil^2$ непобитых белых клеток.

Легко привести пример, показывающий точность нашей оценки: для этого нужно расставить ладьи по одной в каждую строку и каждый столбец, у которых есть граничная черная клетка.

\claim{Ответ} $\lceil n/2\rceil^2$.

\item Пусть $a_1$, $a_2$, \ldots, $a_{99}$ --- произвольная перестановка чисел 1, 2, \ldots, 99. Докажите, что среди чисел $|a_1-1|$, $|a_2-2|$, \ldots, $|a_{99}-99|$ есть два одинаковых.

\claim{Решение} Предположим противное: пусть все числа $|a_1-1|$, $|a_2-2|$, \ldots, $|a_{99}-99|$ различны. Тогда это в точности все числа от 0 до 99. Их сумма равна $99\cdot 49$ --- нечетна. С другой стороны, сумма всех чисел $|a_1-1|$, $|a_2-2|$, \ldots, $|a_{99}-99|$ имеет ту же четность, что и сумма чисел $a_1-1$, $a_2-2$, \ldots, $a_{99}-99$, а последняя сумма равна 0, т.е. четна. Получаем противоречие.

\item В правильном десятиугольнике проведены все диагонали. Возле каждой вершины и возле каждой точки пересечения диагоналей поставлено число $+1$ (рассматриваются только сами диагонали, а не их продолжения). Разрешается одновременно изменить все знаки у чисел, стоящих на одной стороне или на одной диагонали. Можно ли с помощью нескольких таких операций изменить все знаки на противоположные?

\claim{Решение} Обозначим вершины нашего десятиугольника $A_1A_2\ldots A_{10}$. Рассмотрим диагонали $A_1A_4$, $A_2A_6$ и $A_3A_9$ и обозначим через $X$, $Y$ и $Z$ их попарные точки пересечения. Тогда в начальный момент времени у трех точек $X$, $Y$ и $Z$ написаны $+1$. Через точки $X$, $Y$ и $Z$ не проходит больше никакая диагональ, поэтому за каждый шаг мы сможем изменить знак лишь у двух из этих точек. Значит, у точек $X$, $Y$ и $Z$ на каждом шаге будет написано нечетное количество $+1$. Поэтому изменить все знаки на противоположные невозмлжно.

\claim{Ответ} Нет.

\item Алиса и Боб играют в игру. Игровое поле представляет из себя клетчатую полоску размером $1\times 2022$. Игроки по очереди выписывают в пустую клетку любую из букв О и Г. Побеждает тот, после чьего хода в трех соседних клетках появятся буквы ОГО. Если все клетки заполнены, а слова ОГО нет, игра заканчивается вничью. Каков будет исход при правильной игре? 

\claim{Решение} Докажем, что Боб выиграет. Для этого ему нужно создать комбинацию $O\square\square O$. Ясно, что не позднее чем за 2 хода он сможет это сделать, расположив сначала букву О на расстоянии 10 клеток от буквы Алисы, а вторым ходом --- расположив букву О на расстоянии 2 клетки от своей первой буквы О. После этого момента он будет ждать, когда Алиса сделает свой ход в одну из клеток между его буквами. Алиса вынуждена будет это сделать, поскольку после любого ее хода на доске останется нечетное количество пустых клеток, а значит, найдется клетка, справа и слева от которой либо нет букв, либо есть обе буквы. Боб может поставить Г в эту клетку и не проиграет.

\claim{Ответ} Боб.

\item Дана колода карточек. На каждой карточке написано одно из чисел 1, 2, \ldots, $n$. Известно, что сумма всех чисел на всех карточках равна $k\cdot n!$ для некоторого натурального $k$. Докажите, что можно разложить карточки на $k$ стопок так, чтобы сумма чисел, записанных на карточках в каждой стопке, равнялась бы $n!$.

\claim{Решение} Будем доказывать утверждение задачи индукцией по $n$. База очевидна, докажем переход. Как известно, среди любых $n+1$ чиселможно выбрать несколько, сумма которых кратна $n+1$. Выберем такой набор из чисел, написанных на карточках, чтобы их сумма была бы кратна $n+1$ и не превосходила бы $n(n+1)$ (этого всегда можно добиться, т.к. сумма может превосходить $n(n+1)$, только когда на всех карточках написано число $n+1$, а в этом случае достаточно взять одну карточку). Продолжим этот процесс, пока возможно. В результате мы разделим нашу колоду на $s$ кучек, в каждой из которых сумма чисел на карточках равна $\ell_i(n+1)$, где $1\leq\ell_i\leq s$ (в последней группе у нас останется не больше $n$ карточек, каждое из чисел на которых не больше $n+1$). Т.к. $\sum\ell_i (n+1)$ делится на $(n+1)$!, то $\sum\ell_i$ делится на $n!$. Рассмотрим получившиеся кучки как единые карточки с числами $\ell_1$, \ldots, $\ell_s$ и применим предположение индукции. В результате получаем требуемое.

\end{problems}

